
\subsection{$A_K$ 同态的构造 II}
\label{subsec:7-5-5}

在第~\ref{subsec:7-5-4} 小节中，
我们构造了一个 $A_2$ 同态
$(\mathfrak f_1,\mathfrak f_2)$。
本小节继续构造一个 $A_K$ 同态。
先回顾
$(\mathfrak f_1,\ldots,\mathfrak f_K)$
成为从 $(\mathbb R\mathcal X_L,\mathfrak m)$
到 $\Omega(L)$ 的 $A_K$ 同态所需满足的条件。
注意，$\Omega(L)$ 是德拉姆微分分次代数，
因而所有 $\mathfrak m_k$（$k\geq3$）均为零。
这使 $A_K$ 公式得到简化；我们需要证明的公式如下：
\begin{align}
&\sum_{i,\ell}
(-1)^{\deg'P_1+\cdots+\deg'P_i}
\mathfrak f_{k-\ell+1}
\bigl(P_1,\ldots,
 \mathfrak m_\ell(P_{i+1},\ldots,P_{i+\ell}),
 \ldots,P_k\bigr)
\notag\\
&\quad={}
(-1)^\ast d\bigl(\mathfrak f_k(P_1,\ldots,P_k)\bigr)
\notag\\
&\qquad
+\sum_\ell(-1)^{\ast_\ell}
\mathfrak f_\ell(P_1,\ldots,P_\ell)
\wedge
\mathfrak f_{k-\ell}(P_{\ell+1},\ldots,P_k),
\tag{7.5.41}\label{eq:7-5-41}
\end{align}
其中
\[
\ast=\sum_i\deg'P_i,
\qquad
\ast_\ell=
\left(\sum_{i=1}^{\ell}\deg'P_i+1\right)
\left(\sum_{i=\ell+1}^{k}\deg'P_i\right).
\]
注意，左端包含 $\ell=1$ 的项；此时
$\mathfrak m_1=\pm\delta$。

本小节的目标是对 $k$ 作递增归纳、
对 $\sum_i\deg P_i$ 作递减归纳，
构造满足式~\eqref{eq:7-5-41} 的
$\{\mathfrak f_k\}_{k\leq K}$。
式~\eqref{eq:7-5-41} 中含有 $\mathfrak f_k$ 的项，
分别是左端 $\ell=1$ 的情形和右端第一项；
也就是说，它们正是含有 $\delta$ 或 $d$ 的项。

$\mathfrak f_k$ 的构造同第~\ref{subsec:7-5-4} 小节
$\mathfrak f_2$ 的构造相似。
为此需要构造空间 $\mathcal N'_{k+1}$。
回忆
\[
\mathcal N'_2=\{\mathrm{point}\},
\qquad
\mathcal N'_3=[0,1].
\]
在构造 $\mathfrak m_k$ 时，
我们使用了模空间 $\mathcal M_{k+1}^{\mathrm{main}}$。
它的边界是若干
\[
\mathcal M_{\ell+1}^{\mathrm{main}}
\times
\mathcal M_{k-\ell+2}^{\mathrm{main}}
\]
之并。
我们已经证明
$\mathcal M_{k+1}^{\mathrm{main}}$
同胚于 $D^{k-2}$（见第 2.1.2 小节）。
因此，把其边界分解为上述乘积之并，
便给出 $\mathcal M_{k+1}^{\mathrm{main}}$ 的一个胞腔分解。
这个胞腔分解恰同斯塔谢夫为研究 $A_\infty$ 结构
而在 \textbf{[Sta63]} 中引入的斯塔谢夫胞腔分解一致。
斯塔谢夫胞腔边界的每个顶维胞腔，
对应于 $A_\infty$ 公式中不含 $\mathfrak m_1$ 的一项。
在第 7.1.5 小节的定理~\ref{thm:7-1-51} 中，
我们构造了与 $A_\infty$ 映射或同态相关的
$\mathcal N_{k+1}$。

将在第 7.5.6 小节构造的 $\mathcal N'_{k+1}$，
其边界由 $k',k''<k$ 时
$\mathcal N'_{k'+1}$ 与
$\mathcal M_{k''+1}^{\mathrm{main}}$ 的乘积描述。
这同 $\mathcal N_{k+1}$ 的情形相似。
$\mathcal N_{k+1}$ 与 $\mathcal N'_{k+1}$ 之所以不同，
是因为德拉姆复形 $\Omega(L)$ 是微分分次代数，
其所有高阶运算 $\mathfrak m_k$（$k\geq3$）都为零。
正如第 7.5.6 小节的定理 7.5.63 将说明的，
$\mathcal N'_{k+1}$ 可由 $\mathcal N_{k+1}$
压缩其中若干胞腔而得到。

假设 $\ell<k$ 时 $\mathcal N'_{\ell+1}$ 已经定义。
于是 $\partial\mathcal N'_{k+1}$ 可作如下描述。
令
\[
\underline{k}=\{1,\ldots,k\},
\qquad
\underline{k}^{+}=\underline{k}\cup\{\ast\},
\]
其中 $\ast$ 是一个符号。
对每个 $A\in(\underline{k}^{+})^\ell$，
取 $\mathcal N'_{\ell+1}$ 的一个副本并记为
$\mathcal N'_A$；
又固定 $\mathcal M_{\#A+1}^{\mathrm{main}}$
的一个副本并记为 $\mathcal M_A^{\mathrm{main}}$。
那么 $\partial\mathcal N'_{k+1}$
是以下两类胞腔之并：
\begin{align}
&\mathcal M_{(i+1,\ldots,i+\ell)}^{\mathrm{main}}
\times
\mathcal N'_{(1,\ldots,i,\ast,i+\ell+1,\ldots,k)},
\tag{7.5.42.1}\label{eq:7-5-42-1}\\
&\mathcal N'_{(1,\ldots,\ell)}
\times
\mathcal N'_{(\ell+1,\ldots,k)},
\qquad 1\leq\ell\leq k-1.
\tag{7.5.42.2}\label{eq:7-5-42-2}
\end{align}
式~\eqref{eq:7-5-42-1} 对应于
式~\eqref{eq:7-5-41} 中的项
\[
\mathfrak f_{k-\ell+1}
\bigl(P_1,\ldots,
\mathfrak m_\ell(P_{i+1},\ldots,P_{i+\ell}),
\ldots,P_k\bigr),
\]
而式~\eqref{eq:7-5-42-2} 对应于项
\[
\mathfrak f_\ell(P_1,\ldots,P_\ell)
\wedge
\mathfrak f_{k-\ell}(P_{\ell+1},\ldots,P_k).
\]

\begin{proposition}{7.5.43}\label{prop:7-5-43}
存在一个同 $D^{k-1}$（即 $k-1$ 维圆盘）微分同胚的
胞腔复形 $\mathcal N'_{k+1}$，
其边界分解为
式~\eqref{eq:7-5-42-1} 与
式~\eqref{eq:7-5-42-2} 所述胞腔之并。
\end{proposition}

命题~\ref{prop:7-5-43} 等价于说：
$D^{k-1}$ 存在一个光滑胞腔分解，
使式~\eqref{eq:7-5-42-1} 与
式~\eqref{eq:7-5-42-2} 成立。

我们将在第 7.5.6 小节证明命题~\ref{prop:7-5-43}。
在那里将给出 $\mathcal N'_{k+1}$ 的两种定义。
第一种利用 $\mathcal N'_{k+1}$ 与
$\mathcal M_{k+2}^{\mathrm{main}}$ 的关系；
第二种利用 $\mathcal N'_{k+1}$ 与
$\mathcal N_{k+1}$ 的关系，
其中后者已在定理~\ref{thm:7-1-51}、
定义~\ref{def:7-1-54} 中定义。
定理 7.5.63 将说明这两种定义给出同一个空间。
这里先按上述第二种方式描述 $\mathcal N'_{k+1}$。

在 $\mathcal N_{k+1}$ 上定义如下等价关系 $\sim$。
设
\[
\mathbf S=(\Sigma,\vec z,\rho),
\qquad
\mathbf S'=(\Sigma',\vec z',\rho')
\in\mathcal N_{k+1}.
\]
\begin{enumerate}[label=\textup{(7.5.44.\arabic*)},leftmargin=6em]
\item\label{eq:7-5-44-1}
若 $1\notin\operatorname{Im}(\rho)$ 或
$1\notin\operatorname{Im}(\rho')$，则
$\mathbf S\sim\mathbf S'$ 当且仅当
$\mathbf S=\mathbf S'$。
\item\label{eq:7-5-44-2}
设 $1\in\operatorname{Im}(\rho)$ 且
$1\in\operatorname{Im}(\rho')$。
若删去所有时间分配为 $1$ 的分支后二者同构，
则称 $\mathbf S\sim\mathbf S'$。
\end{enumerate}
定义
\[
\mathcal N'_{k+1}=\mathcal N_{k+1}/\!\sim.
\]
第 7.5.6 小节将证明的定理 7.5.63
蕴含 $\mathcal N'_{k+1}$ 是一个胞腔。
由式 (7.1.52) 可知，
$\partial\mathcal N'_{k+1}$ 由
式~\eqref{eq:7-5-42-1}、\eqref{eq:7-5-42-2}
描述
\translatornote{源书把这里的第一处引用印作
“(7.5.42)”，而相邻两处完整列举及命题 7.5.43
都明确给出 (7.5.42.1)、(7.5.42.2) 两类胞腔；
译文补全第一处分项号。见疑误 A-0626。}。

暂且假定这一事实。
本小节据此归纳构造从
$(\mathbb R\mathcal X_L,\mathfrak m)$ 到 $\Omega(L)$ 的
$A_K$ 同态
$(\mathfrak f_1,\ldots,\mathfrak f_K)$。

对 $\vec P=(P_1,\ldots,P_k)$，置
\[
U_{k+1}(\mathfrak f,\vec P)
=\bigl(U_{k+1}\times\mathcal N'_{k+1}\bigr)
\mathbin{\times_{L^{k+1}}}
(P_1\times\cdots\times P_k\times L).
\]
我们将在
\[
\begin{aligned}
&(L\times\mathcal N'_{k+1})
\mathbin{\times_{L^{k+1}}}
(P_1\times\cdots\times P_k\times L)\\
&\hspace{7em}\cong
\mathcal N'_{k+1}\times\bigcap_{i=1}^{k}P_i
\end{aligned}
\]
上构造仓西结构，
使 $U_{k+1}(\mathfrak f,\vec P)$ 成为它的仓西邻域。
障碍丛为
\[
E_{k+1}(\mathfrak f)
=\mathcal N'_{k+1}\times E_{k+1}.
\]

\begin{lemma}{7.5.45}\label{lem:7-5-45}
每个空间
\[
(L\times\mathcal N'_{k+1})
\mathbin{\times_{L^{k+1}}}
(P_1\times\cdots\times P_k\times L)
\]
都有一个仓西结构，且
\[
\bigl(U_{k+1}(\mathfrak f,\vec P),
E_{k+1}(\mathfrak f)\bigr)
\]
是其仓西坐标图。
此外，在边界上，对各种 $k'$ 考察空间
\[
(L\times\mathcal N'_{k'+1})
\mathbin{\times_{L^{k'+1}}}
(P_{i_1}\times\cdots\times P_{i_{k'}}\times L)
\]
的仓西结构，对各种 $k''$ 考察空间
\[
(L\times\mathcal M_{k''+1}^{\mathrm{main}})
\mathbin{\times_{L^{k''+1}}}
(P_{i_1}\times\cdots\times P_{i_{k''}}\times L)
\]
的仓西结构；这些仓西结构的乘积同
\[
(L\times\mathcal N'_{k+1})
\mathbin{\times_{L^{k+1}}}
(P_1\times\cdots\times P_k\times L)
\]
的仓西结构相容。
\end{lemma}

相容性的准确含义以及引理~\ref{lem:7-5-45} 的证明，
都是上一小节 $k=2$ 情形之论证的推广。
这是直接的推广，故略去细节。

现在需要引理~\ref{lem:7-5-37} 的推广，
即下面的引理~\ref{lem:7-5-53}。
先把推广后的陈述说清楚；
一旦有了准确陈述，证明就相当直接。

先注意，空间 $\mathcal N'_{k+1}$ 的分解
由 $\mathcal N_{k+1}$ 的胞腔分解诱导
（见第 7.5.6 小节）。
因此，$\mathcal N'_{k+1}$ 的每个这种胞腔
对应于把若干 $\mathfrak f_\ell$、
$\mathfrak m_\ell$ 与 $\wedge$ 复合的一种方式，
并由此定义映射
\[
(\mathbb R\mathcal X_L)^{\otimes k}\longrightarrow\Omega(L).
\]
例如，当 $k=15$ 时，一种这样的复合是
\begin{align}
&\mathfrak f_4
\bigl(P_1,P_2,
\mathfrak m_3(P_3,P_4,P_5),
\mathfrak m_2(P_6,P_7)\bigr)
\notag\\
&\quad\wedge
\mathfrak f_3
\bigl(P_8,
\mathfrak m_2(\mathfrak m_2(P_9,P_{10}),P_{11}),P_{12}\bigr)
\wedge
\mathfrak f_3(P_{13},P_{14},P_{15}).
\tag{7.5.46}\label{eq:7-5-46}
\end{align}
在微分分次代数 $\Omega(L)$ 中，
我们把 $(A\wedge B)\wedge C$ 与
$A\wedge(B\wedge C)$ 识别，并都写成
$A\wedge B\wedge C$。
这正是我们要在 (7.5.44) 中
压缩时间分配为 $1$ 的分支的理由。
还要指出，式~\eqref{eq:7-5-46} 中不计
$\mathfrak m_1$。

\begin{figure}[!htbp]
 \centering
 \setlength{\unitlength}{0.78mm}
 \begin{picture}(126,105)
  \linethickness{0.45pt}
  \Line(91,52)(116,52)
  \put(116,52){\vector(1,0){7}}
  \Line(91,52)(68,86)
  \Line(91,52)(61,52)
  \Line(91,52)(68,19)
  \Line(68,86)(57,101)
  \Line(68,86)(49,101)
  \Line(68,86)(39,86)
  \Line(68,86)(49,69)
  \Line(39,86)(28,100)
  \Line(39,86)(22,86)
  \Line(39,86)(28,72)
  \Line(49,69)(38,64)
  \Line(49,69)(43,55)
  \Line(61,52)(48,63)
  \Line(61,52)(56,36)
  \Line(61,52)(38,46)
  \Line(38,46)(27,33)
  \Line(38,46)(23,52)
  \Line(23,52)(15,64)
  \Line(23,52)(9,44)
  \Line(68,19)(52,19)
  \Line(68,19)(57,7)
  \Line(68,19)(68,5)
  \put(91,52){\makebox(0,0){$\wedge$}}
  \put(68,86){\makebox(0,0){$F$}}
  \put(39,86){\makebox(0,0){$M$}}
  \put(49,69){\makebox(0,0){$M$}}
  \put(61,52){\makebox(0,0){$F$}}
  \put(38,46){\makebox(0,0){$M$}}
  \put(23,52){\makebox(0,0){$M$}}
  \put(68,19){\makebox(0,0){$F$}}
  \put(59,102){\makebox(0,0)[l]{$P_1$}}
  \put(47,103){\makebox(0,0)[r]{$P_2$}}
  \put(27,102){\makebox(0,0)[r]{$P_3$}}
  \put(20,86){\makebox(0,0)[r]{$P_4$}}
  \put(27,70){\makebox(0,0)[r]{$P_5$}}
  \put(37,64){\makebox(0,0)[r]{$P_6$}}
  \put(42,53){\makebox(0,0)[r]{$P_7$}}
  \put(47,65){\makebox(0,0)[r]{$P_8$}}
  \put(14,66){\makebox(0,0)[r]{$P_9$}}
  \put(8,43){\makebox(0,0)[r]{$P_{10}$}}
  \put(26,31){\makebox(0,0)[r]{$P_{11}$}}
  \put(55,34){\makebox(0,0)[r]{$P_{12}$}}
  \put(50,19){\makebox(0,0)[r]{$P_{13}$}}
  \put(56,5){\makebox(0,0)[r]{$P_{14}$}}
  \put(68,3){\makebox(0,0)[t]{$P_{15}$}}
 \end{picture}
 \par\smallskip
 \renewcommand{\thefigure}{7.5.1}
 \refstepcounter{figure}
 \textbf{图~\thefigure}\label{fig:7-5-1}
\end{figure}

对每个这样的分支，
按下述方式配上组合数据
$([\Sigma],\mathfrak T)$（见图~\ref{fig:7-5-1}）。
设
\[
\Sigma=\bigcup_{\alpha\in I}\Sigma_\alpha
\in\mathcal M_{k+1}^{\mathrm{main}},
\]
其中各 $\Sigma_\alpha$ 都是圆盘。
忘掉 $\Sigma$ 的复结构，
以 $[\Sigma]$ 表示 $\Sigma$ 的组合型（拓扑型）。
按定义~\ref{def:4-6-6}，
在分支指标集 $I$ 上定义偏序。

令
\[
\mathfrak T\colon I\longrightarrow\{M,F,\Lambda\}
\]
把每个分支 $\Sigma_\alpha$ 指派为
$\{M,F,\Lambda\}$ 中的一种类型。
要求 $\mathfrak T$ 满足下列条件。
类型 $M,F,\Lambda$ 分别对应运算
$\mathfrak m,\mathfrak f,\wedge$；
以下条件对应于这些运算可以复合的先后次序。
\begin{enumerate}[label=\textup{(7.5.47.\arabic*)},leftmargin=6em]
\item\label{eq:7-5-47-1}
若 $\Sigma_{\alpha_1}\leq\Sigma_{\alpha_2}$ 且
$\alpha_1\neq\alpha_2$，则
\[
\bigl(\mathfrak T(\alpha_1),\mathfrak T(\alpha_2)\bigr)
\in\{(M,M),(M,F),(M,\Lambda),(F,\Lambda)\}.
\]
\item\label{eq:7-5-47-2}
若 $\mathfrak T(\alpha)=\Lambda$，则
$\alpha=\alpha_0$；这里 $\Sigma_{\alpha_0}$
是含第 $0$ 个标记点的分支。
\item\label{eq:7-5-47-3}
若 $\Sigma_{\alpha_0}$ 是含第 $0$ 个标记点的分支，则
$\mathfrak T(\alpha_0)=\Lambda$ 或
$\mathfrak T(\alpha_0)=F$。
\item\label{eq:7-5-47-4}
若 $\mathfrak T(\alpha_1)=M$ 且
$\mathfrak T(\alpha_2)=\Lambda$，则
$\Sigma_{\alpha_1}\cap\Sigma_{\alpha_2}=\varnothing$。
\end{enumerate}
以 $\mathfrak{CP}_k$ 表示所有这种
$([\Sigma],\mathfrak T)$ 之集。

每个
$([\Sigma],\mathfrak T)\in\mathfrak{CP}_k$
都对应于 $\mathcal N'_{k+1}$ 的一个分支。
例如，对应于式~\eqref{eq:7-5-46} 的
$([\Sigma],\mathfrak T)$ 就是图~\ref{fig:7-5-1}。

一般地，对 $\mathfrak f_*$、$\mathfrak m_*$
与 $\wedge$ 的每一种复合，
可如下配上
$([\Sigma],\mathfrak T)\in\mathfrak{CP}_k$。
$\Sigma$ 的每个分支对应于
$\mathfrak f_*$、$\mathfrak m_*$、$\wedge$
中的一个运算。
注意，把运算
\[
A_1\otimes A_2\otimes\cdots\otimes A_{\ell-1}
\otimes A_\ell
\longmapsto A_1\wedge\cdots\wedge A_\ell
\]
看作一个运算，并为它指派一个分支。
具体规则如下：
\begin{itemize}[leftmargin=2.2em]
\item 若复合中含有 $\wedge$，置
$\mathfrak T(\alpha_0)=\Lambda$；
此时，迭代的 $\wedge$ 运算自动对应于
$\Sigma_{\alpha_0}$。
\item 若 $\alpha$ 所对应的运算是 $\mathfrak m_*$，
置 $\mathfrak T(\alpha)=M$。
\item 若 $\alpha$ 所对应的运算是 $\mathfrak f_*$，
置 $\mathfrak T(\alpha)=F$。
\end{itemize}
容易看出，如此得到的
$([\Sigma],\mathfrak T)$ 满足 (7.5.47)。
反过来，任何满足这些条件的
$([\Sigma],\mathfrak T)$ 都对应于
$\mathcal N'_{k+1}$ 的一个分支。
把这个分支记为
$\mathcal N'_{k+1}([\Sigma],\mathfrak T)$，
并以 $\mathfrak f_{([\Sigma],\mathfrak T)}$
表示相应的
$\mathfrak m_*$、$\mathfrak f_*$ 与 $\wedge$ 的复合。

现在假设每个 $\ell<k$ 都已有 $W_\ell$，
并且对每个 $P_i'\in\mathcal X_L$，
已有仓西邻域
$U_{\ell+1}(\mathfrak f,\vec P')$
上的一个 $W_\ell$ 参数化截面族
$\mathfrak t_{\mathfrak f,\ell+1}^{\vec P'}$。
这里
$\vec P'=(P'_1,\ldots,P'_\ell)$，
而 $U_{\ell+1}(\mathfrak f,\vec P')$
是下列空间上某仓西结构的仓西邻域：
\[
(L\times\mathcal N'_{\ell+1})
\mathbin{\times_{L^{\ell+1}}}
(P'_1\times\cdots\times P'_\ell\times L).
\]
它们如下定义 $\mathfrak f_\ell$。置
\begin{equation}
\mathfrak X_\ell(\vec P')
=\bigl\{(w,p)\in
W_\ell\times U_{\ell+1}(\mathfrak f,\vec P')
\bigm|
\mathfrak t_{\mathfrak f,\ell+1}^{\vec P'}(w,p)=0
\bigr\}.
\tag{7.5.48}\label{eq:7-5-48}
\end{equation}
存在投影
$\pi_{W_\ell}\colon\mathfrak X_\ell(\vec P')\to W_\ell$，
以及由第 $0$ 个标记点处评价映射诱导的映射
$ev\colon\mathfrak X_\ell(\vec P')\to L$。
定义
\begin{equation}
\mathfrak f_\ell(\vec P')
=ev_!\bigl(\pi_{W_\ell}^*(\omega_\ell)\bigr)
\in\Omega(L),
\tag{7.5.49}\label{eq:7-5-49}
\end{equation}
其中 $\omega_\ell$ 是 $W_\ell$ 上的概率测度。

下面定义 $W_k$ 以及
$U_{k+1}(\mathfrak f,\vec P)$ 上的
$W_k$ 参数化截面族
$\mathfrak t_{\mathfrak f,k+1}^{\vec P}$。
先构造参数空间 $W_k$。
思路如下。
运算 $\mathfrak m_*$ 由障碍丛的单个截面定义，
而不是由截面族定义；另一方面，
$\mathfrak f_\ell$ 由 $W_\ell$ 参数化截面族定义。
这意味着，截面族
$\mathfrak t_{\mathfrak f,k+1}^{\vec P}$
在分支
$\mathcal N'_{k+1}([\Sigma],\mathfrak T)$ 上的限制，
由那些满足 $\mathfrak T(\alpha)=F$ 的分支
$\Sigma_\alpha$ 所对应的各个 $W_*$ 之积参数化。
再乘上一个维数足够大的空间 $W_k^{(0)}$，
所得乘积空间就是 $W_k$。
为给出准确的定义，先准备一些记号。

对自然数 $a<b$，置
\[
[a,b]_{\mathbb Z}
=\mathbb Z\cap[a,b]
=\{a,a+1,\ldots,b-1,b\}.
\]
对 $k\in\mathbb Z_{\geq1}$，考察
\begin{equation}
[a,b]_{\mathbb Z}
\subset\underline{k}=\{1,2,\ldots,k\},
\tag{7.5.50.1}\label{eq:7-5-50-1}
\end{equation}
以及它的分解
\begin{equation}
[a,b]_{\mathbb Z}
=\bigcup_{i=1}^{m}[a_i,b_i]_{\mathbb Z},
\tag{7.5.50.2}\label{eq:7-5-50-2}
\end{equation}
并要求
\begin{equation}
[a_i,b_i]_{\mathbb Z}\cap[a_j,b_j]_{\mathbb Z}
=\varnothing
\qquad(i\neq j).
\tag{7.5.50.3}\label{eq:7-5-50-3}
\end{equation}
排列 $a_i$，使得
\begin{equation}
i<j\quad\Longrightarrow\quad a_i<a_j.
\tag{7.5.50.4}\label{eq:7-5-50-4}
\end{equation}
于是 $a_{i+1}=b_i+1$。
允许某些 $i$ 满足 $a_i=b_i$；此时
$[a_i,b_i]_{\mathbb Z}=\{a_i\}$。
不过，要排除 $a=1$、$b=k$ 且所有 $i$ 都满足
$a_i=b_i$ 的情形。换言之，要求
\begin{equation}
a_i\neq b_i\text{ 对某个 }i\text{ 成立},
\qquad\text{或}\qquad
(a,b)\neq(1,k).
\tag{7.5.50.5}\label{eq:7-5-50-5}
\end{equation}

对每个 $k$，考察所有满足
\eqref{eq:7-5-50-1}--\eqref{eq:7-5-50-5} 的
\[
\mathfrak a=
\bigl((a,b);(a_1,b_1),\ldots,(a_m,b_m)\bigr)
\]
之集，并把它记为 $\mathfrak A_k$。
置 $\#\mathfrak a=m$。
对每个 $\mathfrak a\in\mathfrak A_k$，
取 $W_{\#\mathfrak a}$ 的一个副本并记为
$W_{\mathfrak a}$。
定义
\begin{equation}
W_k=W_k^{(0)}\times
\prod_{\mathfrak a\in\mathfrak A_k}W_{\mathfrak a},
\tag{7.5.51}\label{eq:7-5-51}
\end{equation}
其中 $W_k^{(0)}$ 是一个高维空间，稍后再确定。
注意，由条件~\eqref{eq:7-5-50-5}，
$\#\mathfrak a\leq k-1$。
所以式~\eqref{eq:7-5-51} 右端的各个
$W_{\mathfrak a}$ 已由归纳假设定义。

接下来，在 $U_{k+1}(\mathfrak f,\vec P)$ 上定义
$W_k$ 参数化截面族
$\mathfrak t_{\mathfrak f,k+1}^{\vec P}$。
设
$([\Sigma],\mathfrak T)\in\mathfrak{CP}_k$。
对每个满足 $\mathfrak T(\alpha)=F$ 的
$\alpha\in I$，如下定义
$\mathfrak A_k$ 的元素 $\mathfrak a(\alpha)$。

把 $\Sigma_\alpha$ 上的全部奇点或标记点按
$\partial\Sigma_\alpha$ 的循环次序写成
$p_1,\ldots,p_\ell,p_0$。
这里，若 $\alpha\neq\alpha_0$，则 $p_0$ 是满足
\[
p_0\in\Sigma_\alpha\cap\Sigma_{\alpha'},
\qquad
\Sigma_\alpha\leq\Sigma_{\alpha'}
\]
的奇点；若 $\alpha=\alpha_0$，则 $p_0$ 是第 $0$ 个标记点
\translatornote{源书无条件称 $p_0$ 为上述奇点；
但定义 4.6.6 使含第 $0$ 个标记点的分支
$\Sigma_{\alpha_0}$ 成为最大元，因而它没有这样的上方奇点；
同段紧接着还明确讨论光滑单分支 $F$ 型的
$\mathfrak a(\alpha_0)$，此时只能以第 $0$ 个标记点作输出端。
译文据此补出根分支情形。见疑误 A-0627。}。

若 $p_i$ 是奇点，
考察以 $p_i$ 为根的分支树中所含的全部标记点。
定义 $a_i,b_i$，使这个集合的元素恰为第
$a_i,a_i+1,\ldots,b_i$ 个标记点；
注意，其中至少有一个标记点。
若 $p_i$ 是第 $a$ 个标记点，则置
$a_i=b_i=a$。
然后定义
\[
\mathfrak a(\alpha)
=\bigl((a_1,b_\ell);
(a_1,b_1),\ldots,(a_\ell,b_\ell)\bigr).
\]
若 $\Sigma$ 是奇异的，容易看出
$\mathfrak a(\alpha)\in\mathfrak A_k$。
注意，如果 $\Sigma$ 不可约且光滑，
其唯一分支 $\Sigma_0$ 的类型为 $F$，那么
\[
\mathfrak a(\alpha_0)
=\bigl((1,k);\{1\},\ldots,\{k\}\bigr)
\notin\mathfrak A_k.
\]
置
\[
W_{([\Sigma],\mathfrak T)}
=\prod_{\alpha:\,\mathfrak T(\alpha)=F}
W_{\mathfrak a(\alpha)}.
\]

下面的引理~\ref{lem:7-5-53} 由归纳假设推出。
为陈述它，还需一些记号。
令
$\mathfrak t_{([\Sigma],\mathfrak T)}^{\vec P}$
是障碍丛
\[
E_{k+1}(\vec P)
\longrightarrow
U_{k+1}(\mathfrak f,\vec P)([\Sigma],\mathfrak T)
\]
的一个 $W_{([\Sigma],\mathfrak T)}$ 参数化截面族，其中
\[
\begin{aligned}
&U_{k+1}(\mathfrak f,\vec P)([\Sigma],\mathfrak T)\\
&\quad:=
\bigl(U_{k+1}\times
\mathcal N'_{k+1}([\Sigma],\mathfrak T)\bigr)
\mathbin{\times_{L^{k+1}}}
(P_1\times\cdots\times P_k\times L).
\end{aligned}
\]
回忆
$\mathcal N'_{k+1}([\Sigma],\mathfrak T)$
是 $\mathcal N'_{k+1}$ 中对应于
$([\Sigma],\mathfrak T)\in\mathfrak{CP}_k$ 的分支。
置
\begin{align}
&\mathfrak X_k(\vec P)([\Sigma],\mathfrak T)
\notag\\
&\quad={}
\bigl\{(w,p)\in
W_{([\Sigma],\mathfrak T)}\times
U_{k+1}(\mathfrak f,\vec P)([\Sigma],\mathfrak T)
\bigm|
\mathfrak t_{([\Sigma],\mathfrak T)}^{\vec P}
(w,p)=0\bigr\}.
\tag{7.5.52}\label{eq:7-5-52}
\end{align}
存在映射
\[
\pi_{W_{([\Sigma],\mathfrak T)}}\colon
\mathfrak X_k(\vec P)([\Sigma],\mathfrak T)
\longrightarrow W_{([\Sigma],\mathfrak T)},
\qquad
ev\colon
\mathfrak X_k(\vec P)([\Sigma],\mathfrak T)
\longrightarrow L.
\]

\begin{lemma}{7.5.53}\label{lem:7-5-53}
若 $\Sigma$ 是奇异的，则存在障碍丛在
$U_{k+1}(\mathfrak f,\vec P)([\Sigma],\mathfrak T)$
上的一个
$W_{([\Sigma],\mathfrak T)}$ 参数化截面族
$\mathfrak t_{([\Sigma],\mathfrak T)}^{\vec P}$，使得
\[
\mathfrak f_{([\Sigma],\mathfrak T)}(\vec P)
=ev_!\left(
\pi_{W_{([\Sigma],\mathfrak T)}}^*\omega
\right).
\]
这里 $\omega$ 是各因子
$W_{\mathfrak a(\alpha)}$ 上相应概率测度的直积。
\end{lemma}

回忆 $\mathfrak f_{([\Sigma],\mathfrak T)}$
是对应于 $([\Sigma],\mathfrak T)$ 的
$\mathfrak f_*$、$\mathfrak m_*$ 与 $\wedge$ 的复合
\translatornote{源书在引理公式中直接使用未定义的 $\omega$；
按式 (7.5.49) 及
$W_{([\Sigma],\mathfrak T)}$ 的乘积定义，
它只能是各因子上既有概率测度的直积，译文补明于引理中。
紧接的原文又把本节一直使用的 $\mathfrak m_*$
写成 $\overline{\mathfrak m}_*$，
与第 7.5.4 小节开篇的去横线约定及本节全部公式不合，
故统一恢复为 $\mathfrak m_*$。
见疑误 A-0628、A-0629。}。
证明由归纳假设立即得到。

注意，式~\eqref{eq:7-5-51} 中的
$W_{([\Sigma],\mathfrak T)}$ 都是 $W_k$ 的因子。
所以，每个截面族
$\mathfrak t_{([\Sigma],\mathfrak T)}^{\vec P}$
都可视为 $W_k$ 参数化截面族。
把它们胶合起来，得到下列空间上的
$W_k$ 参数化截面族：
\[
\bigl(U_{k+1}\times
\partial\mathcal N'_{k+1}([\Sigma],\mathfrak T)\bigr)
\mathbin{\times_{L^{k+1}}}
(P_1\times\cdots\times P_k\times L).
\]
将其记为
$\mathfrak t_{\mathfrak f,k+1,\partial}^{\vec P}$。

\begin{lemma}{7.5.54}\label{lem:7-5-54}
可以选择 $W_k^{(0)}$，使
$\mathfrak t_{\mathfrak f,k+1,\partial}^{\vec P}$
延拓为一个 $W_k$ 参数化截面族
$\mathfrak t_{\mathfrak f,k+1}^{\vec P}$，
并具有如下横截性。
若置
\begin{equation}
\mathfrak X_k(\vec P)
=\bigl\{(w,p)\in
W_k\times U_{k+1}(\mathfrak f,\vec P)
\bigm|
\mathfrak t_{\mathfrak f,k+1}^{\vec P}(w,p)=0
\bigr\},
\tag{7.5.55}\label{eq:7-5-55}
\end{equation}
则它满足：
\begin{enumerate}[label=\textup{(7.5.56.\arabic*)},leftmargin=6em]
\item\label{eq:7-5-56-1}
$\mathfrak X_k(\vec P)$ 是光滑子流形；
\item\label{eq:7-5-56-2}
$ev_0\colon\mathfrak X_k(\vec P)\to L$ 是浸没。
\end{enumerate}
\end{lemma}

其证明是引理~\ref{lem:7-5-35} 与
引理~\ref{lem:7-5-37} 的证明的直接推广。

\begin{definition}{7.5.57}\label{def:7-5-57}
定义
\[
\mathfrak f_k(\vec P)
=ev_!\bigl(\pi_{W_k}^*(\omega_k)\bigr)
\in\Omega(L),
\]
其中
\[
ev\colon\mathfrak X_k(\vec P)\to L,
\qquad
\pi_{W_k}\colon\mathfrak X_k(\vec P)\to W_k,
\]
而 $\omega_k$ 是 $W_k$ 上次数为
$\dim W_k$ 的紧支撑微分形式，且满足
\[
\int_{W_k}\omega_k=1.
\]
\end{definition}

\begin{lemma}{7.5.58}\label{lem:7-5-58}
同态 $\mathfrak f_k$ 满足式~\eqref{eq:7-5-41}。
\end{lemma}

这立即由引理~\ref{lem:7-5-10}、
\ref{lem:7-5-53}、\ref{lem:7-5-45} 推出。
这样，对每个固定的 $K$，
我们都构造出了一个 $A_K$ 同态。
再重复第 7.2.9--7.2.11 小节的论证，
便完成定理~\ref{thm:7-5-1} 的证明。

利用第 7.2.10--7.2.13 小节的论证，
还可证明本节构造的同伦等价在相差同伦的意义下
不依赖于我们所作的各种选择。
